%%%%%%%%%%%%%%%%%%%%%%%%%%%%%%%%%%%%%%%%%%%%%%%%%%%%%%%%%%%%%%%%%%%%%%%%%%%%%%%
%%%%%%%%%%%%%%%%%%%%%%%%%%%%%%%%%%%%%%%%%%%%%%%%%%%%%%%%%%%%%%%%%%%%%%%%%%%%%%%
% APPENDIX
%%%%%%%%%%%%%%%%%%%%%%%%%%%%%%%%%%%%%%%%%%%%%%%%%%%%%%%%%%%%%%%%%%%%%%%%%%%%%%%
%%%%%%%%%%%%%%%%%%%%%%%%%%%%%%%%%%%%%%%%%%%%%%%%%%%%%%%%%%%%%%%%%%%%%%%%%%%%%%%
\newpage
\appendix
% \onecolumn

\section{Hyperparameters}

We used 50 or 10 samples from PG19 for most experiments.
Across all settings, the base language model was frozen; only the compression-token parameters were optimized (either the compression embeddings, or the low-dimensional embedding together with its projection).
We optimized with AdamW. For each model, we swept the learning rate over \(\{0.01, 0.1, 0.5, 1.0\}\); the resulting metrics are summarized in Table~\ref{tab:all_learning_rates}. For the low-dimensional embedding and projection variant, we used a learning rate of 0.01.
We ran up to 10{,}000 optimization steps per sample in all experiments; for progressive runs, we additionally capped the optimization of each newly added token at 1{,}000 steps.
We used 100 warmup steps for all models.

\section{Compute}

Each experiment ran on a single NVIDIA A100 80GB GPU, with up to 16 GPUs used in parallel across experiments. Training time depends on the sequence length and loss type, ranging from a few minutes (short sequences, small models) to approximately 24 hours (long sequences, large models). The total compute used for this research (including debugging and unsuccessful experiments) is approximately 2800 GPU-hours.


\section{Full cramming vs progressive.}
\label{app:1}

In Table~\ref{tab:full_vs_progressive_appendix} we present results from Section~\ref{sec:experiments} for full and compressive cramming comparison for different model sizes of Llama and Pythia families.

\begin{table}[!htbp]
    \caption{Full vs. progressive cramming on PG19. ``Full'' cramming uses a fixed token budget and may stop below perfect reconstruction \cite{kuratov2025cramming}; progressive cramming fixes reconstruction at 100\% and reports the achieved token count.}
    \label{tab:full_vs_progressive_appendix}
    \begin{center}
    \begin{small}
    % paper-lint: n_samples=10
\begin{tabular}{llrr}
\toprule
 & & \multicolumn{2}{c}{Accuracy (\%)} \\
\cmidrule(lr){3-4}
 Type                                           & Tokens                  & TF                         & Greedy                     \\
\midrule
 \multicolumn{4}{l}{\textbf{Llama-3.2-1B}} \\
 Full                                           & 512                     & 99.80 {\small $\pm$ 0.00}  & 1.51 {\small $\pm$ 0.94}   \\
 Prog.                                          & 402 {\small $\pm$ 85}   & 100.00                     & 100.00                     \\
 \midrule
 \multicolumn{4}{l}{\textbf{Llama-3.2-3B}} \\
 Full                                           & 1024                    & 99.60 {\small $\pm$ 0.70}  & 1.02 {\small $\pm$ 1.23}   \\
 Prog.                                          & 902 {\small $\pm$ 207}  & 100.00                     & 100.00                     \\
 \midrule
 \multicolumn{4}{l}{\textbf{Llama-3.1-8B}} \\
 Full                                           & 1568                    & 99.72 {\small $\pm$ 0.32}  & 1.29 {\small $\pm$ 1.12}   \\
 Prog.                                          & 1064 {\small $\pm$ 394} & 100.00                     & 100.00                     \\
 \midrule \midrule
 \multicolumn{4}{l}{\textbf{Pythia160m}} \\
 Full                                           & 32                      & 68.44 {\small $\pm$ 17.54} & 26.39 {\small $\pm$ 24.08} \\
 Prog.                                          & 11 {\small $\pm$ 2}     & 100.00                     & 100.00                     \\
 \midrule
 \multicolumn{4}{l}{\textbf{Pythia410m}} \\
 Full                                           & 128                     & 89.06 {\small $\pm$ 6.15}  & 13.95 {\small $\pm$ 23.93} \\
 Prog.                                          & 102 {\small $\pm$ 41}   & 100.00                     & 100.00                     \\
 \midrule
 \multicolumn{4}{l}{\textbf{Pythia1.4b}} \\
 Full                                           & 256                     & 98.48 {\small $\pm$ 3.06}  & 39.23 {\small $\pm$ 43.79} \\
 Prog.                                          & 544 {\small $\pm$ 57}   & 100.00                     & 100.00                     \\
\bottomrule
\end{tabular}

    \end{small}
    \end{center}
\vskip -0.1in
\end{table}

\begin{table*}[tb]
    \caption{Results of the learning-rate sweep across models and settings. ``Compressed Tokens'' is the achieved prefix length $n$, ``Trajectory Length'' is $L_{\text{traj}}$, and ``PCA 99\%'' is the number of principal components explaining 99\% of trajectory variance.}
    \label{tab:all_learning_rates}
    \begin{center}
    \begin{small}
    \begin{tabular}{lllll}
\toprule
 Model                         & Compressed Tokens           & Information Gain         & Trajectory Length         & PCA 99\%                    \\
\midrule
 Llama-3.1-8B                  & 1298.4 {\small $\pm$ 537.5} & 3760 {\small $\pm$ 1418} & 738 {\small $\pm$ 265}    & 40.8 {\small $\pm$ 13.16}  \\
 Llama-3.1-8B {\small lr=0.1}  & 1063.5 {\small $\pm$ 394.4} & 3028 {\small $\pm$ 1321} & 4861 {\small $\pm$ 1033}  & 74.4 {\small $\pm$ 7.94}   \\
 Llama-3.1-8B {\small lr=0.5}  & 934.4 {\small $\pm$ 123.2}  & 2758 {\small $\pm$ 486}  & 15424 {\small $\pm$ 1437} & 138 {\small $\pm$ 11.78}   \\
 Llama-3.1-8B {\small lr=1.0}  & 1286.4 {\small $\pm$ 524.8} & 3381 {\small $\pm$ 1521} & 29061 {\small $\pm$ 4605} & 186.6 {\small $\pm$ 13.59} \\
 \midrule
 pythia-1.4b                   & 188.2 {\small $\pm$ 27.8}   & 581 {\small $\pm$ 101}   & 83 {\small $\pm$ 14}      & 15.7 {\small $\pm$ 1.85}   \\
 pythia-1.4b {\small lr=0.1}   & 274.3 {\small $\pm$ 56.2}   & 846 {\small $\pm$ 106}   & 956 {\small $\pm$ 123}    & 23.2 {\small $\pm$ 1.99}   \\
 pythia-1.4b {\small lr=0.5}   & 543.8 {\small $\pm$ 56.8}   & 1694 {\small $\pm$ 125}  & 6730 {\small $\pm$ 390}   & 49 {\small $\pm$ 3.32}     \\
 pythia-1.4b {\small lr=1.0}   & 558.2 {\small $\pm$ 74.1}   & 1737 {\small $\pm$ 158}  & 11381 {\small $\pm$ 993}  & 63.3 {\small $\pm$ 6.21}   \\
 \midrule
 SmolLM2-1.7B {\small lr=0.1}  & 370.2 {\small $\pm$ 113.1}  & 1119 {\small $\pm$ 350}  & 1027 {\small $\pm$ 167}   & 29.2 {\small $\pm$ 2.93}   \\
 SmolLM2-1.7B {\small lr=0.5}  & 540.5 {\small $\pm$ 139.4}  & 1450 {\small $\pm$ 350}  & 5225 {\small $\pm$ 827}   & 66.8 {\small $\pm$ 9.14}   \\
 SmolLM2-1.7B {\small lr=1.0}  & 502.5 {\small $\pm$ 180.6}  & 1428 {\small $\pm$ 473}  & 9442 {\small $\pm$ 1810}  & 83 {\small $\pm$ 16.34}    \\
 gemma-3-4b-pt {\small lr=0.1} & 286.5 {\small $\pm$ 164.6}  & 949 {\small $\pm$ 466}   & 1044 {\small $\pm$ 328}   & 31.8 {\small $\pm$ 9.21}   \\
 gemma-3-4b-pt {\small lr=0.5} & 498.6 {\small $\pm$ 86.5}   & 1437 {\small $\pm$ 308}  & 6055 {\small $\pm$ 925}   & 71.6 {\small $\pm$ 11.34}  \\
 \midrule
 gemma-3-4b-pt {\small lr=1.0} & 558 {\small $\pm$ 86.1}     & 1258 {\small $\pm$ 511}  & 12173 {\small $\pm$ 1004} & 101.6 {\small $\pm$ 10.66} \\
\bottomrule
\end{tabular}

    \end{small}
    \end{center}
  \vskip -0.1in
\end{table*}

\section{Progressive across models scales}

Table~\ref{tab:progressive_for_model_scales} shows that larger models tend to use more trajectory directions (higher PCA 99\%).


\begin{table*}[tb]
    \caption{Progressive cramming trajectory statistics on PG19. ``Compressed Tokens'' is the achieved prefix length $n$, ``Trajectory Length'' is $L_{\text{traj}}$, and ``PCA 99\%'' is the number of principal components explaining 99\% of trajectory variance.}
    \label{tab:progressive_for_model_scales}
    \begin{center}
    \begin{small}
    \begin{tabular}{lllll}
\hline
 Model                         & Compressed Tokens           & Information Gain         & Trajectory Length        & PCA 99\%                   \\
\hline
 Llama-3.2-1B {\small lr=0.1}  & 402.2 {\small $\pm$ 84.8}   & 1500 {\small $\pm$ 282}  & 1736 {\small $\pm$ 206}  & 44.1 {\small $\pm$ 5.73}  \\
 Llama-3.2-3B {\small lr=0.1}  & 902.2 {\small $\pm$ 206.8}  & 3074 {\small $\pm$ 527}  & 4232 {\small $\pm$ 913}  & 61.5 {\small $\pm$ 4.54}  \\
 Llama-3.1-8B {\small lr=0.1}  & 1063.5 {\small $\pm$ 394.4} & 3028 {\small $\pm$ 1321} & 4861 {\small $\pm$ 1033} & 74.4 {\small $\pm$ 7.94}  \\
 \midrule
 pythia-160m {\small lr=0.5}   & 10.7 {\small $\pm$ 2}       & 19 {\small $\pm$ 26}     & 652 {\small $\pm$ 196}   & 5.1 {\small $\pm$ 1.64}   \\
 pythia-410m {\small lr=0.5}   & 102.1 {\small $\pm$ 40.8}   & 323 {\small $\pm$ 105}   & 3613 {\small $\pm$ 841}  & 37.5 {\small $\pm$ 11.16} \\
 pythia-1.4b {\small lr=0.5}   & 543.8 {\small $\pm$ 56.8}   & 1694 {\small $\pm$ 125}  & 6730 {\small $\pm$ 390}  & 49 {\small $\pm$ 3.32}    \\
 \midrule
 SmolLM2-135M {\small lr=0.1}  & 38.5 {\small $\pm$ 14.1}    & 168 {\small $\pm$ 66}    & 178 {\small $\pm$ 40}    & 12.2 {\small $\pm$ 2.96}  \\
 SmolLM2-360M {\small lr=0.1}  & 61 {\small $\pm$ 24.6}      & 266 {\small $\pm$ 122}   & 234 {\small $\pm$ 111}   & 13.4 {\small $\pm$ 3.95}  \\
 SmolLM2-1.7B {\small lr=0.1}  & 370.2 {\small $\pm$ 113.1}  & 1119 {\small $\pm$ 350}  & 1027 {\small $\pm$ 167}  & 29.2 {\small $\pm$ 2.93}  \\
 \midrule
 gemma-3-270m {\small lr=0.1}  & 71.2 {\small $\pm$ 14.2}    & 392 {\small $\pm$ 71}    & 259 {\small $\pm$ 32}    & 18.5 {\small $\pm$ 2.54}  \\
 gemma-3-1b-pt {\small lr=0.1} & 74.8 {\small $\pm$ 34.3}    & 338 {\small $\pm$ 104}   & 386 {\small $\pm$ 89}    & 18.9 {\small $\pm$ 6.59}  \\
 gemma-3-4b-pt {\small lr=0.1} & 286.5 {\small $\pm$ 164.6}  & 949 {\small $\pm$ 466}   & 1044 {\small $\pm$ 328}  & 31.8 {\small $\pm$ 9.21}  \\
\hline
\end{tabular}

    \end{small}
    \end{center}
    \vskip -0.1in

\end{table*}

\section{Prefix tuning}

Prefix tuning introduces a separate trainable prefix embedding at each transformer layer, avoiding the per-layer forward-pass bottleneck of soft-prompt setups such as full or progressive cramming. The trade-off is parameter cost: it requires $L$ times more trainable embeddings, where $L$ is the number of hidden layers in the LLM.

To separate ``cramming'' from general soft-prompt steering, we evaluate prefix tuning baselines.
Table~\ref{tab:prefix_tuning_attention_hijacking} reports the same attention-mass metric for prefix tokens.
Compared to cramming, prefix tuning typically attracts less attention mass than BOS and shows weaker or more variable correlation patterns, indicating that attention hijacking is not an inevitable consequence of prepending learned embeddings but depends on the optimization objective.

We also evaluated compression limits under prefix tuning. Table~\ref{tab:prefix_tuning_accuracy} shows that, for Llama~3 models at 3B and 8B, we cannot reach the cramming limits; however, in the prefix-tuning setup these models can still compress up to 8k tokens with perfect reconstruction accuracy, while 16k tokens results in out-of-memory (OOM) failures. For the 1B Llama~3 model, the limit is about 4096 tokens, which is roughly \(10\times\) higher than with progressive cramming. This gain comes at a higher parameter cost: prefix tuning requires \(L\)-times more trainable embeddings (one per layer), i.e., \(16\times\) more parameters for a 16-layer model.
For Pythia, the prefix-tuning setup increases the compression limit by about \(20\times\) for Pythia-410M and by about \(8\times\) for Pythia-1.4B relative to progressive cramming.


\begin{table*}[tb]
    \centering
    \caption{Prefix-tuning attention mass, computed with the same metric as Table~\ref{tab:attn_hijacking} but for learned prefix tokens instead of the compression embedding.}
    \label{tab:prefix_tuning_attention_hijacking}
    \begin{small}
    \begin{tabular}{llllr}
\toprule
 Model        & Compression Token (\%)     & BOS Token Original (\%)    & Diff (\%)                   &   Correlation \\
\midrule
 Llama-3.2-3B & 23.24 {\small $\pm$ 4.55} & 67.75 {\small $\pm$ 1.18} & -44.51 {\small $\pm$ 5.61} &        0.5641 \\
 SmolLM2-135M & 26.53 {\small $\pm$ 2.78} & 41.74 {\small $\pm$ 6.35} & -15.21 {\small $\pm$ 7.13} &        0.8691 \\
 SmolLM2-360M & 28.07 {\small $\pm$ 4.8}  & 49.81 {\small $\pm$ 7.67} & -21.74 {\small $\pm$ 8.13} &        0.6066 \\
 SmolLM2-1.7B & 36.43 {\small $\pm$ 4.58} & 44.03 {\small $\pm$ 6.73} & -7.61 {\small $\pm$ 8.38}  &        0.9102 \\
 Qwen3-4B     & 1.95 {\small $\pm$ 0.22}  & 50.84 {\small $\pm$ 1.54} & -48.89 {\small $\pm$ 1.46} &        0.0771 \\
\bottomrule
\end{tabular}

    \end{small}
\vskip -0.1in
\end{table*}


\begin{table}[tb]
    \centering
    \caption{Prefix-tuning reconstruction accuracy compared to progressive cramming. ``Tokens'' denotes the number of learned prefix tokens (or reconstructed tokens for progressive cramming).}
    \label{tab:prefix_tuning_accuracy}
    \begin{small}
    \begin{tabular}{llll}
\hline
 Experiment      & Type          & Tokens                  & Accuracy                   \\
\hline
 \midrule
 Llama-3.2-1B         & Progr.        & 402 {\small $\pm$ 85}   & 1.0                        \\
 Llama-3.2-1B         & Full PrefixT. & 4096                    & 1 {\small $\pm$ 0}         \\
 Llama-3.2-1B         & Full PrefixT. & 8192                    & 0.977 {\small $\pm$ 0.032} \\
 Llama-3.2-1B         & Full PrefixT. & 16384                   & 0.767 {\small $\pm$ 0.105} \\
 \midrule
 Llama-3.2-3B         & Progr.        & 902 {\small $\pm$ 207}  & 1.0                        \\
 Llama-3.2-3B         & Full PrefixT. & 8192                    & 1 {\small $\pm$ 0}         \\
 \midrule
 Llama-3.1-8B         & Progr.        & 1064 {\small $\pm$ 394} & 1.0                        \\
 Llama-3.1-8B         & Full PrefixT. & 8192                    & 1 {\small $\pm$ 0}         \\
 \midrule
 Pythia160m          & Progr.        & 11 {\small $\pm$ 2}     & 1.0                        \\
 Pythia160m          & Full PrefixT. & 1024                    & 0.488 {\small $\pm$ 0.06}  \\
 Pythia160m          & Full PrefixT. & 2048                    & 0.394 {\small $\pm$ 0.034} \\
 Pythia160m          & Full PrefixT. & 4096                    & 0.263 {\small $\pm$ 0.025} \\
 Pythia160m          & Full PrefixT. & 8192                    & 0.168 {\small $\pm$ 0.018} \\
 Pythia160m          & Full PrefixT. & 16384                   & 0.139 {\small $\pm$ 0.026} \\
 \midrule
 Pythia410m          & Progr.        & 102 {\small $\pm$ 41}   & 1.0                        \\
 Pythia410m          & Full PrefixT. & 2048                    & 1 {\small $\pm$ 0}         \\
 Pythia410m          & Full PrefixT. & 4096                    & 0.964 {\small $\pm$ 0.071} \\
 Pythia410m          & Full PrefixT. & 8192                    & 0.588 {\small $\pm$ 0.106} \\
 Pythia410m          & Full PrefixT. & 16384                   & 0.418 {\small $\pm$ 0.158} \\
 \midrule
 Pythia1.4b          & Progr.        & 544 {\small $\pm$ 57}   & 1.0                        \\
 Pythia1.4b          & Full PrefixT. & 4096                    & 1 {\small $\pm$ 0}         \\
 Pythia1.4b          & Full PrefixT. & 8192                    & 0.979 {\small $\pm$ 0.015} \\
 Pythia1.4b          & Full PrefixT. & 16384                   & 0.663 {\small $\pm$ 0.12}  \\
\hline
\end{tabular}

    \end{small}
\vskip -0.1in
\end{table}





\section{Dataset Modifications}

To probe sensitivity to surface form, we create controlled variants of PG19 samples and compare their progressive trajectories (Figure~\ref{fig:trajectories_text_modifications}).
Each variant shares an identical 64-token prefix and modifies the suffix in one of three ways: random word permutation, greedy decoding from Llama-3.1-8B (``Sampled''), or lowercasing.
Trajectories coincide during the shared prefix and diverge immediately after the edit point, indicating that the optimization path is highly sensitive to token-level details.
Notably, lowercasing (which should preserve high-level semantics) induces a markedly different trajectory, suggesting that progressive cramming is driven by tokenization- and distribution-level properties rather than semantic equivalence.



\begin{figure*}[tb]
    \centering
    \centerline{\includegraphics[width=\linewidth]{Llama3.1-8B-text-modifications_lr-0p1.pdf}}
    \caption{Progressive trajectories projected onto PCA components (Llama-3.1-8B, one PG19 sample). All variants share the same 64-token prefix; the suffix is either original (Base), randomly permuted (Random), greedily sampled continuation (Sampled), or lowercased (Lowercased). Circles mark the start and squares mark the end. Axis titles report the signed explained-variance fraction.}
    \label{fig:trajectories_text_modifications}
\end{figure*}



% \begin{table}[]
%     \centering
%     \caption{Progressive cramming on modified PG19 samples (Llama-3.1-8B). All variants share a 64-token prefix and differ only in the suffix (Base/Random/Sampled/Lowercased). {\color{red} TODO recompute via \texttt{scripts/paper/text\_modifications\_table.py}.}}
%     \label{tab:text_modifications}

% \begin{tabular}{llll}
% \hline
%  Experiment   & Tokens                  & Info Gain                & PCA 99\%              \\
% \hline
%  Base         & 1048 {\small $\pm$ 404} & 3169 {\small $\pm$ 1212} & 74 {\small $\pm$ 8}  \\
%  Random       & 1 {\small $\pm$ 0}      & 0 {\small $\pm$ 0}       & nan                  \\
%  Sampled      & 2048 {\small $\pm$ 0}   & 432 {\small $\pm$ 141}   & 60 {\small $\pm$ 28} \\
%  Lowercased   & 1 {\small $\pm$ 0}      & 0 {\small $\pm$ 0}       & nan                  \\
% \hline
% \end{tabular}
% \end{table}



\section{Full activation alignment experiments}

Table~\ref{tab:full_activation_alignment_and_low_dim_projections} reports the full sweep over activation alignment strength \(\alpha\), number of aligned layers \(L\), and projection dimension, together with the resulting progressive cramming metrics.

\begin{table*}[tb]
    \centering
    \caption{Full sweep of activation alignment and low-dimensional projection hyperparameters. ``$\alpha$'' is alignment strength, ``$L$'' is the number of aligned layers, and ``dim'' is projection dimension.}\label{tab:full_activation_alignment_and_low_dim_projections}

    \begin{small}
    \begin{tabular}{lllll}
\hline
 Model                                                               & Compressed Tokens           & Information Gain         & Trajectory Length         & PCA 99\%                  \\
\hline
 Llama-3.1-8B {\small $\alpha=1.0$} {\small $L=2$}                   & 1532.3 {\small $\pm$ 280}   & 4694 {\small $\pm$ 341}  & 812 {\small $\pm$ 167}    & 45.8 {\small $\pm$ 5.29} \\
 Llama-3.1-8B {\small $\alpha=1.0$} {\small $L=4$}                   & 1593.6 {\small $\pm$ 306.6} & 4961 {\small $\pm$ 603}  & 958 {\small $\pm$ 207}    & 42.9 {\small $\pm$ 3.01} \\
 Llama-3.1-8B {\small $\alpha=1.0$} {\small $L=8$}                   & 1633.3 {\small $\pm$ 275.4} & 5130 {\small $\pm$ 247}  & 1146 {\small $\pm$ 137}   & 40 {\small $\pm$ 1.26}   \\
 Llama-3.1-8B {\small $\alpha=1.0$} {\small $L=16$}                  & 1487.3 {\small $\pm$ 255.4} & 4686 {\small $\pm$ 276}  & 1146 {\small $\pm$ 112}   & 33.6 {\small $\pm$ 2.2}  \\
 Llama-3.1-8B {\small $\alpha=1.0$} {\small $L=24$}                  & 664.8 {\small $\pm$ 99.2}   & 2095 {\small $\pm$ 247}  & 570 {\small $\pm$ 87}     & 20.5 {\small $\pm$ 1.75} \\
 Llama-3.1-8B {\small $\alpha=1.0$} {\small $L=32$}                  & 192.7 {\small $\pm$ 50.6}   & 614 {\small $\pm$ 169}   & 150 {\small $\pm$ 31}     & 12.6 {\small $\pm$ 1.69} \\
 Llama-3.1-8B {\small dim=32} {\small $\alpha=1.0$} {\small $L=8$}   & 1804.8 {\small $\pm$ 297.8} & 5103 {\small $\pm$ 748}  & 13584 {\small $\pm$ 3823} & 34.7 {\small $\pm$ 5.37} \\
 Llama-3.1-8B {\small dim=64} {\small $\alpha=1.0$} {\small $L=8$}   & 1853.3 {\small $\pm$ 256.2} & 5512 {\small $\pm$ 456}  & 14367 {\small $\pm$ 3209} & 30.8 {\small $\pm$ 5.27} \\
 Llama-3.1-8B {\small dim=128} {\small $\alpha=1.0$} {\small $L=8$}  & 1858.9 {\small $\pm$ 283.6} & 5145 {\small $\pm$ 1695} & 18967 {\small $\pm$ 3766} & 33.5 {\small $\pm$ 2.84} \\
 Llama-3.1-8B {\small dim=256} {\small $\alpha=1.0$} {\small $L=8$}  & 1810.6 {\small $\pm$ 236}   & 5105 {\small $\pm$ 1693} & 23576 {\small $\pm$ 5844} & 34.1 {\small $\pm$ 5.5}  \\
 \midrule
 pythia-1.4b {\small $\alpha=1.0$} {\small $L=4$}                    & 182.7 {\small $\pm$ 31.1}   & 553 {\small $\pm$ 63}    & 83 {\small $\pm$ 9}       & 14.6 {\small $\pm$ 1.96} \\
 pythia-1.4b {\small $\alpha=1.0$} {\small $L=8$}                    & 178.7 {\small $\pm$ 33.1}   & 524 {\small $\pm$ 92}    & 80 {\small $\pm$ 13}      & 14.3 {\small $\pm$ 1.79} \\
 pythia-1.4b {\small $\alpha=1.0$} {\small $L=16$}                   & 155.5 {\small $\pm$ 35.5}   & 427 {\small $\pm$ 66}    & 69 {\small $\pm$ 16}      & 13.5 {\small $\pm$ 1.5}  \\
 pythia-1.4b {\small $\alpha=1.0$} {\small $L=20$}                   & 116.6 {\small $\pm$ 26.1}   & 307 {\small $\pm$ 35}    & 58 {\small $\pm$ 10}      & 11.9 {\small $\pm$ 0.83} \\
 pythia-1.4b {\small dim=32} {\small $\alpha=1.0$} {\small $L=8$}    & 423.1 {\small $\pm$ 104.4}  & 1317 {\small $\pm$ 221}  & 2498 {\small $\pm$ 1438}  & 15.5 {\small $\pm$ 4.57} \\
 pythia-1.4b {\small dim=64} {\small $\alpha=1.0$} {\small $L=8$}    & 447.8 {\small $\pm$ 81.9}   & 1398 {\small $\pm$ 186}  & 3309 {\small $\pm$ 2569}  & 16.9 {\small $\pm$ 4.01} \\
 pythia-1.4b {\small dim=128} {\small $\alpha=1.0$} {\small $L=8$}   & 468.3 {\small $\pm$ 68.6}   & 1448 {\small $\pm$ 106}  & 4699 {\small $\pm$ 2414}  & 20 {\small $\pm$ 4.96}   \\
 pythia-1.4b {\small dim=256} {\small $\alpha=1.0$} {\small $L=8$}   & 498.2 {\small $\pm$ 92.7}   & 1555 {\small $\pm$ 146}  & 5396 {\small $\pm$ 3309}  & 19.6 {\small $\pm$ 5.89} \\
 \midrule
 SmolLM2-1.7B {\small $\alpha=1.0$} {\small $L=8$}                   & 148.1 {\small $\pm$ 56.7}   & 432 {\small $\pm$ 157}   & 75 {\small $\pm$ 18}      & 11.9 {\small $\pm$ 1.97} \\
 SmolLM2-1.7B {\small $\alpha=1.0$} {\small $L=4$}                   & 186.9 {\small $\pm$ 47.3}   & 530 {\small $\pm$ 104}   & 85 {\small $\pm$ 8}       & 13.3 {\small $\pm$ 1.73} \\
 SmolLM2-1.7B {\small $\alpha=1.0$} {\small $L=16$}                  & 126.5 {\small $\pm$ 63.3}   & 366 {\small $\pm$ 155}   & 66 {\small $\pm$ 22}      & 11.3 {\small $\pm$ 2.19} \\
 SmolLM2-1.7B {\small $\alpha=1.0$} {\small $L=20$}                  & 96.1 {\small $\pm$ 56.2}    & 254 {\small $\pm$ 114}   & 50 {\small $\pm$ 20}      & 8.8 {\small $\pm$ 3.54}  \\
 SmolLM2-1.7B {\small dim=32} {\small $\alpha=1.0$} {\small $L=8$}   & 292.7 {\small $\pm$ 113.3}  & 851 {\small $\pm$ 327}   & 1144 {\small $\pm$ 705}   & 11.1 {\small $\pm$ 3.99} \\
 SmolLM2-1.7B {\small dim=64} {\small $\alpha=1.0$} {\small $L=8$}   & 321.4 {\small $\pm$ 148.3}  & 914 {\small $\pm$ 305}   & 1668 {\small $\pm$ 2014}  & 10.4 {\small $\pm$ 2.29} \\
 SmolLM2-1.7B {\small dim=128} {\small $\alpha=1.0$} {\small $L=8$}  & 363.2 {\small $\pm$ 67.2}   & 1065 {\small $\pm$ 213}  & 2236 {\small $\pm$ 1198}  & 10.4 {\small $\pm$ 1.5}  \\
 SmolLM2-1.7B {\small dim=256} {\small $\alpha=1.0$} {\small $L=8$}  & 468 {\small $\pm$ 243.7}    & 1350 {\small $\pm$ 621}  & 6544 {\small $\pm$ 9765}  & 11.6 {\small $\pm$ 3.61} \\
 \midrule
 gemma-3-4b-pt {\small $\alpha=1.0$} {\small $L=4$}                  & 60.3 {\small $\pm$ 52.7}    & 224 {\small $\pm$ 163}   & 60 {\small $\pm$ 20}      & 9.6 {\small $\pm$ 3.01}  \\
 gemma-3-4b-pt {\small $\alpha=1.0$} {\small $L=8$}                  & 75.1 {\small $\pm$ 84.9}    & 268 {\small $\pm$ 256}   & 65 {\small $\pm$ 44}      & 9.3 {\small $\pm$ 3.93}  \\
 gemma-3-4b-pt {\small $\alpha=1.0$} {\small $L=16$}                 & 55.8 {\small $\pm$ 42.8}    & 217 {\small $\pm$ 156}   & 57 {\small $\pm$ 19}      & 8.4 {\small $\pm$ 2.65}  \\
 gemma-3-4b-pt {\small $\alpha=1.0$} {\small $L=20$}                 & 48.3 {\small $\pm$ 31.4}    & 191 {\small $\pm$ 104}   & 55 {\small $\pm$ 16}      & 8.3 {\small $\pm$ 2.45}  \\
 gemma-3-4b-pt {\small dim=32} {\small $\alpha=1.0$} {\small $L=8$}  & 121.3 {\small $\pm$ 127.9}  & 393 {\small $\pm$ 317}   & 429 {\small $\pm$ 277}    & 9.3 {\small $\pm$ 3.93}  \\
 gemma-3-4b-pt {\small dim=64} {\small $\alpha=1.0$} {\small $L=8$}  & 327.5 {\small $\pm$ 224.5}  & 1082 {\small $\pm$ 708}  & 807 {\small $\pm$ 502}    & 13.4 {\small $\pm$ 3.58} \\
 gemma-3-4b-pt {\small dim=128} {\small $\alpha=1.0$} {\small $L=8$} & 361.3 {\small $\pm$ 214.4}  & 1209 {\small $\pm$ 672}  & 1511 {\small $\pm$ 801}   & 14.5 {\small $\pm$ 2.06} \\
 gemma-3-4b-pt {\small dim=256} {\small $\alpha=1.0$} {\small $L=8$} & 217.1 {\small $\pm$ 146.2}  & 719 {\small $\pm$ 425}   & 1091 {\small $\pm$ 557}   & 13.3 {\small $\pm$ 2.69} \\
 \midrule
 Qwen3-4B {\small $\alpha=1.0$} {\small $L=4$}                       & 166.9 {\small $\pm$ 56.1}   & 699 {\small $\pm$ 193}   & 104 {\small $\pm$ 18}     & 13.8 {\small $\pm$ 1.54} \\
 Qwen3-4B {\small $\alpha=1.0$} {\small $L=8$}                       & 197.1 {\small $\pm$ 74}     & 792 {\small $\pm$ 243}   & 118 {\small $\pm$ 25}     & 15.1 {\small $\pm$ 2.12} \\
 Qwen3-4B {\small $\alpha=1.0$} {\small $L=16$}                      & 170.9 {\small $\pm$ 71.3}   & 711 {\small $\pm$ 265}   & 109 {\small $\pm$ 28}     & 13.9 {\small $\pm$ 2.47} \\
 Qwen3-4B {\small $\alpha=1.0$} {\small $L=20$}                      & 168.9 {\small $\pm$ 39.6}   & 695 {\small $\pm$ 132}   & 106 {\small $\pm$ 12}     & 13.9 {\small $\pm$ 1.37} \\
 Qwen3-4B {\small dim=32} {\small $\alpha=1.0$} {\small $L=8$}       & 774.9 {\small $\pm$ 156.4}  & 3063 {\small $\pm$ 556}  & 4243 {\small $\pm$ 1543}  & 25.9 {\small $\pm$ 4.48} \\
 Qwen3-4B {\small dim=64} {\small $\alpha=1.0$} {\small $L=8$}       & 811.8 {\small $\pm$ 189.1}  & 3181 {\small $\pm$ 567}  & 7120 {\small $\pm$ 2953}  & 30.7 {\small $\pm$ 4.17} \\
 Qwen3-4B {\small dim=128} {\small $\alpha=1.0$} {\small $L=8$}      & 870.8 {\small $\pm$ 160.5}  & 3400 {\small $\pm$ 300}  & 9040 {\small $\pm$ 3358}  & 31 {\small $\pm$ 5.42}   \\
 Qwen3-4B {\small dim=256} {\small $\alpha=1.0$} {\small $L=8$}      & 895.7 {\small $\pm$ 219.4}  & 3470 {\small $\pm$ 476}  & 9813 {\small $\pm$ 3019}  & 32.5 {\small $\pm$ 4.25} \\
\hline
\end{tabular}

    \end{small}
\vskip -0.1in

\end{table*}


\section{Full Low dimensional experiments}

Table~\ref{tab:low_dim_projection_results} ablates projection dimension across model families.
Three takeaways stand out:
\begin{itemize}
    \item For Llama-3.1-8B and SmolLM2-1.7B, projection can substantially increase achievable token counts and information gain compared to the unconstrained baseline.
    \item Increasing projection dimension does not reliably increase PCA 99\% (and can even decrease it), but it often increases the measured trajectory length, suggesting longer optimization paths within a similarly low-dimensional subspace.
    \item For Pythia-1.4B, small projection dimensions strongly reduce PCA 99\% and trajectory length, but also reduce capacity; adding a larger projection partially recovers information gain.
\end{itemize}


\begin{table*}[tb]
    \centering
    \caption{Effect of low-dimensional projection size on progressive cramming. ``dim'' in the model column denotes the projection dimension $k$ (Section~\ref{sec:lowdim}).}
    \label{tab:low_dim_projection_results}

    \begin{small}
    \begin{tabular}{lllll}
\hline
 Model                          & Compressed Tokens           & Information Gain         & Trajectory Length          & PCA 99\%                  \\
\hline
 Llama-3.1-8B {\small lr=0.1}   & 1063.5 {\small $\pm$ 394.4} & 3028 {\small $\pm$ 1321} & 4861 {\small $\pm$ 1033}   & 74.4 {\small $\pm$ 7.94} \\
 Llama-3.1-8B {\small dim=32}   & 1745 {\small $\pm$ 306}     & 5312 {\small $\pm$ 330}  & 10250 {\small $\pm$ 3256}  & 35.7 {\small $\pm$ 5.1}  \\
 Llama-3.1-8B {\small dim=64}   & 1834.7 {\small $\pm$ 251}   & 5694 {\small $\pm$ 333}  & 9977 {\small $\pm$ 2641}   & 32.5 {\small $\pm$ 4.3}  \\
 Llama-3.1-8B {\small dim=128}  & 1834.6 {\small $\pm$ 305.6} & 5538 {\small $\pm$ 537}  & 18713 {\small $\pm$ 10220} & 32.8 {\small $\pm$ 4.6}  \\
 Llama-3.1-8B {\small dim=256}  & 1730.8 {\small $\pm$ 384.3} & 4870 {\small $\pm$ 1680} & 21448 {\small $\pm$ 5583}  & 29.4 {\small $\pm$ 3.29} \\
 Llama-3.1-8B {\small dim=512}  & 1652.9 {\small $\pm$ 404.5} & 5072 {\small $\pm$ 799}  & 25852 {\small $\pm$ 7528}  & 36.2 {\small $\pm$ 5.88} \\
 \midrule
 pythia-1.4b {\small lr=0.5}    & 543.8 {\small $\pm$ 56.8}   & 1694 {\small $\pm$ 125}  & 6730 {\small $\pm$ 390}    & 49 {\small $\pm$ 3.32}   \\
 pythia-1.4b {\small dim=32}    & 358.2 {\small $\pm$ 80.6}   & 1137 {\small $\pm$ 241}  & 1551 {\small $\pm$ 954}    & 15.1 {\small $\pm$ 3.01} \\
 pythia-1.4b {\small dim=64}    & 392.8 {\small $\pm$ 96.6}   & 1240 {\small $\pm$ 286}  & 1920 {\small $\pm$ 1191}   & 16.7 {\small $\pm$ 3.69} \\
 pythia-1.4b {\small dim=128}   & 373 {\small $\pm$ 81.8}     & 1168 {\small $\pm$ 240}  & 1915 {\small $\pm$ 1133}   & 15.8 {\small $\pm$ 3.22} \\
 pythia-1.4b {\small dim=256}   & 375.5 {\small $\pm$ 100.7}  & 1189 {\small $\pm$ 314}  & 2135 {\small $\pm$ 1312}   & 16.5 {\small $\pm$ 2.16} \\
 pythia-1.4b {\small dim=512}   & 415.5 {\small $\pm$ 76.3}   & 1333 {\small $\pm$ 200}  & 2863 {\small $\pm$ 1433}   & 18 {\small $\pm$ 2.49}   \\
 \midrule
 SmolLM2-1.7B {\small lr=0.1}   & 370.2 {\small $\pm$ 113.1}  & 1119 {\small $\pm$ 350}  & 1027 {\small $\pm$ 167}    & 29.2 {\small $\pm$ 2.93} \\
 SmolLM2-1.7B {\small dim=32}   & 335.8 {\small $\pm$ 78.7}   & 1007 {\small $\pm$ 284}  & 543 {\small $\pm$ 218}     & 14.2 {\small $\pm$ 1.99} \\
 SmolLM2-1.7B {\small dim=64}   & 403.2 {\small $\pm$ 77.7}   & 1195 {\small $\pm$ 177}  & 778 {\small $\pm$ 167}     & 13.8 {\small $\pm$ 1.94} \\
 SmolLM2-1.7B {\small dim=128}  & 431.9 {\small $\pm$ 115.6}  & 1252 {\small $\pm$ 288}  & 1236 {\small $\pm$ 653}    & 13.9 {\small $\pm$ 2.51} \\
 SmolLM2-1.7B {\small dim=256}  & 487.7 {\small $\pm$ 152.7}  & 1449 {\small $\pm$ 402}  & 2265 {\small $\pm$ 1405}   & 16 {\small $\pm$ 4.47}   \\
 SmolLM2-1.7B {\small dim=512}  & 557.1 {\small $\pm$ 158.8}  & 1615 {\small $\pm$ 350}  & 2590 {\small $\pm$ 1163}   & 17.7 {\small $\pm$ 3.95} \\
 \midrule
 gemma-3-4b-pt {\small lr=0.1}  & 286.5 {\small $\pm$ 164.6}  & 949 {\small $\pm$ 466}   & 1044 {\small $\pm$ 328}    & 31.8 {\small $\pm$ 9.21} \\
 gemma-3-4b-pt {\small dim=32}  & 233.7 {\small $\pm$ 202.3}  & 670 {\small $\pm$ 428}   & 673 {\small $\pm$ 343}     & 12 {\small $\pm$ 3.1}    \\
 gemma-3-4b-pt {\small dim=64}  & 214.9 {\small $\pm$ 172.8}  & 718 {\small $\pm$ 568}   & 700 {\small $\pm$ 486}     & 11.8 {\small $\pm$ 2.71} \\
 gemma-3-4b-pt {\small dim=128} & 275.5 {\small $\pm$ 222.6}  & 860 {\small $\pm$ 667}   & 1107 {\small $\pm$ 1265}   & 13.3 {\small $\pm$ 2.76} \\
 gemma-3-4b-pt {\small dim=256} & 252.7 {\small $\pm$ 189.9}  & 829 {\small $\pm$ 635}   & 1372 {\small $\pm$ 886}    & 13.1 {\small $\pm$ 4.83} \\
 gemma-3-4b-pt {\small dim=512} & 217.4 {\small $\pm$ 169.1}  & 639 {\small $\pm$ 439}   & 1297 {\small $\pm$ 975}    & 12.3 {\small $\pm$ 5.95} \\
\hline
\end{tabular}

    \end{small}
\vskip -0.1in

\end{table*}

\begin{table*}[t]
\centering
\caption{Compression reconstruction accuracy under non-ideal training convergence.
We report mean teacher-forcing accuracy over all compressed samples, mean greedy
decoding accuracy when generating autoregressively from the compression token,
and the fraction of samples with a token mismatch at indices 0, 1, and 2.
Variants include ignoring the BOS loss during training (no BOS) and upweighting
the loss of the first two tokens by a factor of three (2 leading).}
\label{tab:compression_reconstruction_summary}
\begin{small}
% paper-lint: n_samples=50
\begin{tabular}{lll r c ccc ccc}
\toprule
Model & Tokens & Setup & TF conv. & Greedy conv. & \multicolumn{3}{c}{Greedy mismatch (\%)} & \multicolumn{3}{c}{TF mismatch (\%)} \\
\cmidrule(lr){6-8} \cmidrule(lr){9-11}
 &  &  &  &  & @0 & @1 & @2 & @0 & @1 & @2 \\
\midrule
Llama-3.2-1B & 512 & common & 0.9904 & 0.0019 & 86.0\% & 100.0\% & 100.0\% & 88.0\% & 100.0\% & 28.0\% \\
Llama-3.2-1B & 512 & no BOS & 0.9887 & 0.0380 & 56.0\% & 72.0\% & 82.0\% & 56.0\% & 34.0\% & 16.0\% \\
Llama-3.2-1B & 512 & 2 leading & 0.9903 & 0.0045 & 0.0\% & 100.0\% & 100.0\% & 0.0\% & 100.0\% & 30.0\% \\
\midrule
Llama-3.2-3B & 1024 & common & 0.9902 & 0.0019 & 100.0\% & 100.0\% & 98.0\% & 100.0\% & 98.0\% & 36.0\% \\
Llama-3.2-3B & 1024 & no BOS & 0.9821 & 0.0150 & 80.0\% & 85.0\% & 90.0\% & 80.0\% & 62.0\% & 48.0\% \\
Llama-3.2-3B & 1024 & 2 leading & 0.9904 & 0.0123 & 10.0\% & 98.0\% & 100.0\% & 10.0\% & 98.0\% & 42.0\% \\
\midrule
Llama-3.1-8B & 1568 & common & 0.9907 & 0.0041 & 100.0\% & 100.0\% & 100.0\% & 100.0\% & 100.0\% & 60.0\% \\
Llama-3.1-8B & 1568 & no BOS & 0.9798 & 0.0044 & 84.0\% & 92.0\% & 96.0\% & 84.0\% & 68.0\% & 46.0\% \\
Llama-3.1-8B & 1568 & 2 leading & 0.9908 & 0.0152 & 28.0\% & 96.0\% & 98.0\% & 28.0\% & 96.0\% & 54.0\% \\
\bottomrule
\end{tabular}

\end{small}
\vskip -0.1in
\end{table*}


\section{Analysis of the Recostruction Accuracy}
\label{app:I}

As noted in Section~\ref{sec:prior_limitations}, a 99\% token-level accuracy threshold is insufficient to guarantee correct reconstruction when evaluation is performed without \emph{teacher forcing} (i.e., when generated tokens are fed back into the prefix rather than replacing them with the ground-truth tokens). Table~\ref{tab:compression_reconstruction_summary} shows that under non-progressive training, the remaining $\sim$1\% error is concentrated in the first few positions (indices 0-2); these early mistakes then compound, leading to divergent generations.


\section{Effect of Imperfect Reconstruction on HellaSwag and ARC}
\label{app:imperfect_reconstruction}

The downstream evaluation in Section~\ref{sec:semantic_evaluation} retains benchmark instances whose prefixes are not perfectly reconstructed, which raises the question of whether the observed capability drop is caused by the compression embedding itself disrupting reasoning, or simply by failed reconstruction corrupting the prefix.
To disentangle the effect of imperfect reconstruction from the effect of the compression embedding itself, we emphasize the difference between autoregressive generation and multiple-choice benchmarks such as HellaSwag and ARC. In greedy generation, even small reconstruction errors can lead to large deviations in the output due to error accumulation, especially when errors occur early in the sequence (Section~\ref{sec:prior_limitations}).

In contrast, HellaSwag and ARC operate in a non-generative (scoring) regime: the model evaluates fixed context--continuation pairs without regenerating the prefix. As a result, these benchmarks are significantly less sensitive to small reconstruction imperfections, since no sequential error propagation occurs.

The near-identical token-normalized accuracies for SmolLM2-1.7B (Table~\ref{tab:rebuttal_token_norm}) indicate that imperfect reconstruction has minimal impact on benchmark performance. This confirms that the observed degradation is driven primarily by the compression embedding itself, rather than by reconstruction errors.

\begin{table}[t]
\centering
\caption{Token-normalized accuracy (\%) for HellaSwag, ARC-Easy, and ARC-Challenge on SmolLM2-1.7B. Values are reported for perfectly reconstructed examples and for all examples.}
\label{tab:rebuttal_token_norm}
\begin{small}
\begin{tabular}{lccc}
\toprule
\textbf{Subset} & \textbf{HellaSwag} & \textbf{ARC-Easy} & \textbf{ARC-Challenge} \\
\midrule
\textit{Perfect only} & 36.05 & 55.39 & 28.86 \\
\textit{All samples}  & 36.14 & 55.21 & 28.44 \\
\bottomrule
\end{tabular}

\end{small}
\end{table}

